\chapter{The planet inside the account}\label{ch:planet}

A physical account becomes useful when its numbers can be traced to observations and interpreted within their scope. Planetary assessments offer a demanding example: their quantities differ, their consequences unfold over different timescales, and their projections depend on future choices. The same discipline will be needed when a small EBU model moves from an exact equation to a claim about the world.

\section{A changing inheritance}
The IPCC's 2023 synthesis concluded that human activities had unequivocally caused global warming, with global surface temperature in 2011--2020 about $1.1\,^\circ\mathrm{C}$ above 1850--1900. It assessed widespread changes and unequal contributions and exposure.\cite{intro-ipcc} This is a comparison between specified periods. It should not be presented as today's instantaneous reading.

The IPBES global assessment estimated that roughly one million animal and plant species were threatened with extinction. Threatened does not mean already extinct, and the estimate is not a census of every species.\cite{intro-ipbes} The International Resource Panel's 2024 outlook described a conditional path under which global extraction could rise by almost sixty percent from 2020 to 2060 without urgent concerted action.\cite{intro-resources} The projection names a possible course under assumptions, not an inevitable future.

Temperature, biodiversity and extraction are different objects. Their headline numbers cannot be added to obtain one planetary score. They nevertheless show why an economic action can affect conditions that support people long after the transaction closes.

An EBU account must begin with the same respect for distinctions. It needs a declared physical state and potential, appropriate units and a boundary within which the comparison makes sense. A proof about that representation is valuable, but it cannot establish that the representation includes every important harm.

\section{The short return and the long consequence}
A company can receive revenue from a water withdrawal this month while consequences for groundwater, soils or downstream ecology develop over years. Legal duties and public scrutiny may already address some of those effects. Others may remain uncertain or outside the immediate decision.

Several clocks are involved. A transaction closes at one time. Material moves through a system at another. Regeneration, degradation and ecological response may have still different timescales. A complete receipt for the first event is not evidence that the later processes have ended.

This gives the future EBU system two related jobs. It should state what transition has actually been evaluated, and it should distinguish that record from predictions about later response. A model of pending effects or feedback may help with the prediction. It must supply its own assumptions and evidence; a potential difference does not secretly contain a law of motion.

The benefit could be better-timed maintenance and intervention. A clear distinction between completed improvement, expected recovery and continuing exposure can reduce premature declarations of success. For society, that can matter when a service must remain dependable throughout a recovery, rather than merely look better at its final inspection.

\section{A projection is a conditional warning}
The IPCC compares outcomes under different emissions pathways. A pathway states assumptions about future activity; its appearance in a table does not itself assign it a probability.\cite{intro-ipcc} Comparing alternatives is useful because a decision can still affect which conditions develop.

Our later coalition tables serve a narrower mathematical purpose: they compare specified combinations of actions under one protocol. Those alternatives must not be mistaken for observed histories or probabilities merely because the table is complete. Nor does a model-generated endpoint become a measurement by passing through exact algebra.

This separation supports a practical workflow. Keep observations, response assumptions and valuation rules identifiable. Then a corrected observation can be distinguished from a changed forecast or a newly chosen field. Decisions can improve without silently rewriting what an older record meant.

\begin{workedbox}{Read the year and the condition}
The UNEP/IRP projection of almost sixty percent more extraction by 2060 uses 2020 as its comparison and assumes a continuation without urgent concerted action. It describes neither an already observed 2060 nor an equal increase in every resource.
\end{workedbox}

\section{What is missing from the dashboard?}
A district can report rising income while its aquifer falls, its clinic loses reliable service and its waste accumulates. Each report may be accurate. The error would be to expect one to substitute for all the others.

The constructive question is what information a particular decision requires. Which stocks or qualities can the action change? Which source, receiver and sink belong in its physical boundary? Which effects are already present, which remain pending, and which require a separate forecast? These questions can produce a tractable local account even when a complete model of the planet is unavailable.

Later we shall prove that a local calculation can equal a declared global potential change when unchanged factors cancel. That is an efficiency result with exact assumptions, not permission to ignore a remote consequence that genuinely changes. It offers a way to make a physical account usable without demanding global knowledge from every person opening a valve.

Before developing it, we look at familiar physical examples of local relations with consistent consequences. They will introduce the difference between a descriptive principle, a quantity being compared and a law that actually moves a system.
